\documentclass[11pt,a4paper]{article}

\usepackage[utf8]{inputenc}
\usepackage[T1]{fontenc}
\usepackage[spanish,es-tabla]{babel}
\usepackage[a4paper,margin=2.3cm]{geometry}
\usepackage{fancyhdr}
\usepackage{xcolor}
\usepackage{listings}
\usepackage{booktabs}
\usepackage{array}
\usepackage{longtable}
\usepackage{enumitem}
\usepackage{graphicx}
\usepackage[hidelinks,colorlinks=true,linkcolor=blue!60!black,urlcolor=blue!60!black,citecolor=blue!60!black]{hyperref}
\usepackage{titlesec}

% Colores para listings
\definecolor{codegray}{rgb}{0.95,0.95,0.95}
\definecolor{codeblue}{rgb}{0.13,0.29,0.53}
\definecolor{codegreen}{rgb}{0.0,0.5,0.0}
\definecolor{codeorange}{rgb}{0.8,0.4,0.0}

\lstdefinestyle{shell}{
    backgroundcolor=\color{codegray},
    basicstyle=\ttfamily\small,
    breakatwhitespace=false,
    breaklines=true,
    captionpos=b,
    keepspaces=true,
    showspaces=false,
    showstringspaces=false,
    showtabs=false,
    tabsize=2,
    frame=single,
    framerule=0pt,
    xleftmargin=8pt,
    xrightmargin=8pt,
    aboveskip=8pt,
    belowskip=8pt,
}

\lstdefinelanguage{yaml}{
    keywords={apiVersion,kind,metadata,spec,name,labels,app,replicas,selector,
              matchLabels,template,containers,image,imagePullPolicy,ports,
              containerPort,env,value,type,protocol,port,targetPort,nodePort,services,
              build,context,dockerfile,environment,restart,args,depends_on},
    keywordstyle=\color{codeblue}\bfseries,
    sensitive=true,
    comment=[l]{\#},
    commentstyle=\color{codegreen}\itshape,
    morestring=[b]",
    stringstyle=\color{codeorange},
}

\lstdefinelanguage{bash}{
    morekeywords={sudo,kubectl,k3s,ctr,images,export,import,scale,rollout,
                  cordon,drain,uncordon,deployment,backend,frontend,flatpak,run,
                  org.apache.jmeter,curl,docker,compose,up,down,build,push,
                  scp,ssh,sshpass,grep,awk,sed,tail,head,for,do,done,if,then,fi,
                  echo,sleep,export},
    keywordstyle=\color{codeblue}\bfseries,
    sensitive=true,
    comment=[l]{\#},
    commentstyle=\color{codegreen}\itshape,
    morestring=[b]",
    morestring=[b]',
    stringstyle=\color{codeorange},
}

\lstset{style=shell}

\titleformat{\section}{\Large\bfseries\color{blue!50!black}}{\thesection}{1em}{}
\titleformat{\subsection}{\large\bfseries\color{blue!50!black}}{\thesubsection}{1em}{}

\pagestyle{fancy}
\fancyhf{}
\fancyhead[L]{E-Converse | Análisis de Rendimiento K8s}
\fancyhead[R]{\thepage}
\renewcommand{\headrulewidth}{0.4pt}
\renewcommand{\footrulewidth}{0pt}

\setlist[itemize]{itemsep=2pt,topsep=4pt}
\setlist[enumerate]{itemsep=2pt,topsep=4pt}

\begin{document}

\begin{titlepage}
    \centering
    \vspace*{1.5cm}
    {\huge\bfseries Análisis del Comportamiento de \\[0.3cm] E-Converse bajo Diferentes\\[0.3cm] Escenarios de Despliegue\par}
    \vspace{1cm}
    {\Large Proyecto Final\par}
    \vspace{0.4cm}
    {\large Ingeniería de Software III\par}
    \vspace{2cm}
    {\large Universidad Industrial de Santander\par}
    \vspace{1.5cm}
    \begin{tabular}{l}
        \large\textbf{Equipo:} \\[0.2cm]
        Juan Sebastián Otero --- 2220053 \\
        Daniel Santiago Convers --- 2221120 \\
        Juan David Paipa --- 2220062 \\
        Jhon Anderson Vargas --- 2220086 \\
    \end{tabular}
    \vfill
    {\large Mayo 2026\par}
\end{titlepage}

\tableofcontents
\newpage

% ======================================================
\section{Introducción}
% ======================================================

Este documento es el reporte de entrega del proyecto final del curso. Su propósito es servir como guía operativa: explica qué es la aplicación, cómo está organizado el repositorio, cómo levantar el sistema localmente y en Kubernetes, cómo reproducir las pruebas de rendimiento con JMeter y dónde está el análisis completo de los resultados.

El \textbf{análisis técnico detallado}, las \textbf{tablas} con todas las métricas y las \textbf{gráficas comparativas} se encuentran en el notebook Jupyter \texttt{Analisis\_Rendimiento.ipynb}. Este PDF se enfoca en las instrucciones de ejecución y los enlaces necesarios.

% ======================================================
\section{Descripción de la Aplicación: E-Converse}
% ======================================================

\textbf{E-Converse} es una plataforma de comercio electrónico con backend Spring Boot, frontend React y persistencia en MongoDB Atlas.

\begin{center}
\begin{tabular}{ll}
\toprule
\textbf{Componente} & \textbf{Tecnología} \\
\midrule
Backend & Spring Boot 3.5.6 (Java 20), API REST, JWT \\
Frontend & React 19 + Vite 7, TailwindCSS, Zustand \\
Base de Datos & MongoDB Atlas (10,000+ productos) \\
Contenedorización & Docker (multi-stage) \\
Orquestación & Kubernetes (K3s v1.35.5) \\
Pruebas de Carga & Apache JMeter 5.6.3 \\
Análisis & Jupyter Notebook (pandas, matplotlib, seaborn) \\
\bottomrule
\end{tabular}
\end{center}

\subsection{Funcionalidades principales}

\begin{itemize}
    \item Autenticación con JWT (login, registro, roles administrador/cliente).
    \item Panel administrativo con CRUDs de usuarios, roles, categorías, productos y pedidos.
    \item Catálogo público con filtros por categoría, precio y nombre.
    \item Carrito de compras persistente en MongoDB con índice único por usuario.
    \item Proceso de checkout con simulación de pasarela de pago.
    \item Historial de pedidos por usuario.
\end{itemize}

\subsection{Endpoints utilizados para las pruebas de carga}

\begin{center}
\begin{tabular}{lll}
\toprule
\textbf{Endpoint} & \textbf{Tamaño de respuesta} & \textbf{Característica} \\
\midrule
\texttt{GET /api/categorias/list} & $\sim$800 bytes & Endpoint ligero (4 documentos) \\
\texttt{GET /api/productos/list} & $\sim$3.5 MB & Endpoint pesado (10K documentos) \\
\bottomrule
\end{tabular}
\end{center}

% ======================================================
\section{Estructura del Repositorio}
% ======================================================

\begin{lstlisting}[language=bash]
Proyecto-Spring-Boot-E-converse_vers2/
|-- backend/                    # API Spring Boot (codigo + Dockerfile)
|-- frontend-react/             # SPA React + Vite (codigo + nginx + Dockerfile)
|-- docker-compose.yml          # Despliegue local (Fase 1)
|-- k8s/
|   |-- backend-deployment.yaml
|   |-- frontend-deployment.yaml
|-- jmeter/
|   |-- econverse_load_test.jmx   # Plan JMeter parametrizado
|   |-- run_all_tests.sh          # Script de automatizacion 9 escenarios
|   |-- run_1node_only.sh         # Re-ejecucion solo Fase 2
|   |-- results/                  # 9 CSVs (corrida final, 0% errores)
|   |-- results_old/              # 10 CSVs (corrida inicial, 67% errores)
|-- scripts/
|   |-- populate_db.js            # Insercion de 10K productos
|   |-- setup_roles_admin.js      # Configuracion inicial
|   |-- redeploy-backend.sh       # Re-deploy backend en K3s
|-- Analisis_Rendimiento.ipynb    # Notebook de analisis (PRINCIPAL)
|-- REPORTE_PROYECTO.md           # Diario de problemas y soluciones
|-- Reporte_Final.tex / .pdf      # Este documento
|-- README.md
\end{lstlisting}

% ======================================================
\section{Cómo Ejecutar el Proyecto}
% ======================================================

\subsection{Requisitos previos}

\begin{itemize}
    \item Java 17+ y Maven (o usar el contenedor Docker que ya incluye Maven).
    \item Node.js 18+ y npm (para frontend dev).
    \item Docker y, opcionalmente, Docker Compose v2.
    \item Acceso a un clúster Kubernetes (K3s, MicroK8s, Kind o cloud).
    \item Una URI de MongoDB Atlas o instancia local.
\end{itemize}

\subsection{Opción A: Ejecución local sin contenedores (desarrollo)}

\begin{lstlisting}[language=bash,caption={Backend Spring Boot}]
cd backend
export SPRING_DATA_MONGODB_URI="mongodb+srv://<user>:<pass>@<cluster>/e_converse"
./mvnw spring-boot:run
# Backend disponible en http://localhost:8080
\end{lstlisting}

\begin{lstlisting}[language=bash,caption={Frontend React}]
cd frontend-react
npm install
npm run dev
# Frontend disponible en http://localhost:5174
\end{lstlisting}

\subsection{Opción B: Docker Compose (Fase 1 del proyecto)}

\begin{lstlisting}[language=bash]
export SPRING_DATA_MONGODB_URI="mongodb+srv://<user>:<pass>@<cluster>/e_converse"
docker compose up --build
# Frontend en http://localhost:8081
# Backend  en http://localhost:8080
\end{lstlisting}

\subsection{Opción C: Kubernetes (Fases 2--4)}

\textbf{1. Construir las imágenes localmente:}
\begin{lstlisting}[language=bash]
docker build -t e-converse-backend:latest ./backend
docker build -t e-converse-frontend:latest ./frontend-react
\end{lstlisting}

\textbf{2. Distribuir las imágenes a todos los nodos del clúster} (importante: \emph{todos} los nodos, no solo los workers):
\begin{lstlisting}[language=bash]
docker save e-converse-backend:latest  -o backend.tar
docker save e-converse-frontend:latest -o frontend.tar
# Copiar los tar a cada nodo y luego en cada uno:
sudo k3s ctr images import backend.tar
sudo k3s ctr images import frontend.tar
\end{lstlisting}

\textbf{3. Configurar la URI de MongoDB} en \texttt{k8s/backend-deployment.yaml} (variable \texttt{SPRING\_DATA\_MONGODB\_URI}).

\textbf{4. Desplegar:}
\begin{lstlisting}[language=bash]
sudo k3s kubectl apply -f k8s/backend-deployment.yaml
sudo k3s kubectl apply -f k8s/frontend-deployment.yaml
sudo k3s kubectl get pods -o wide
# La aplicacion queda en http://<IP-master>:30080
\end{lstlisting}

\subsection{Poblar la base de datos con 10K productos}

\begin{lstlisting}[language=bash]
cd scripts
npm install
node setup_roles_admin.js   # Roles y usuario admin
node populate_db.js         # Insercion de 10K productos mock
\end{lstlisting}

% ======================================================
\section{Cómo Visualizar el Notebook}
% ======================================================

El análisis completo, las tablas con todas las métricas y las gráficas comparativas están en \texttt{Analisis\_Rendimiento.ipynb}. Hay tres formas de abrirlo:

\subsection{Opción 1: Jupyter Lab/Notebook local}

\begin{lstlisting}[language=bash]
pip install jupyter pandas matplotlib seaborn
cd /ruta/al/repositorio
jupyter notebook Analisis_Rendimiento.ipynb
# O bien:
jupyter lab
\end{lstlisting}

\subsection{Opción 2: VS Code con extensión Jupyter}

Abrir el archivo \texttt{Analisis\_Rendimiento.ipynb} desde VS Code; la extensión Jupyter detectará el kernel de Python automáticamente.

\subsection{Opción 3: Google Colab}

Subir el notebook a Google Drive, abrir con Colab y subir manualmente la carpeta \texttt{jmeter/results/} a la sesión (o montar Drive). El notebook lee los CSVs desde la ruta relativa \texttt{jmeter/results/}.

\subsection{Re-ejecutar todas las celdas}

Para regenerar todas las gráficas a partir de los CSVs:

\begin{lstlisting}[language=bash]
jupyter nbconvert --to notebook --execute --inplace Analisis_Rendimiento.ipynb
\end{lstlisting}

% ======================================================
\section{Cómo Reproducir las Pruebas de Carga}
% ======================================================

\subsection{Plan de pruebas JMeter (parametrizado)}

El archivo \texttt{jmeter/econverse\_load\_test.jmx} es un plan de pruebas con parámetros configurables desde la línea de comandos:

\begin{center}
\begin{tabular}{lll}
\toprule
\textbf{Parámetro} & \textbf{Flag} & \textbf{Valor por defecto} \\
\midrule
Servidor              & \texttt{-Jserver}   & \texttt{10.6.101.163} \\
Puerto                & \texttt{-Jport}     & \texttt{30080} \\
Hilos concurrentes    & \texttt{-Jthreads}  & \texttt{20} \\
Ramp-up               & \texttt{-Jrampup}   & \texttt{20 s} \\
Duración              & \texttt{-Jduration} & \texttt{60 s} \\
\bottomrule
\end{tabular}
\end{center}

\subsection{Ejecutar una prueba puntual desde la línea de comandos}

\begin{lstlisting}[language=bash]
flatpak run org.apache.jmeter -n \
  -t jmeter/econverse_load_test.jmx \
  -l jmeter/results/mi_prueba.csv \
  -Jserver=10.6.101.163 -Jport=30080 \
  -Jthreads=20 -Jrampup=20 -Jduration=60
\end{lstlisting}

\textit{Si JMeter está instalado nativamente, usar \texttt{jmeter} en lugar de \texttt{flatpak run org.apache.jmeter}.}

\subsection{Ejecutar la suite completa (9 escenarios)}

\begin{lstlisting}[language=bash]
./jmeter/run_all_tests.sh
# Duracion aproximada: 40-50 minutos
# Resultados en jmeter/results/k8s_*.csv
\end{lstlisting}

El script automatiza todo el ciclo: para cada combinación de nodos $\times$ réplicas, hace \texttt{cordon}/\texttt{drain} de los nodos correspondientes, escala el deployment, espera un warmup activo (HTTP 200 consecutivos $+$ 60 s de buffer), corre JMeter y guarda el CSV. Al final restaura el clúster a su estado original.

\subsection{Configuración previa requerida (importante)}

Antes de correr la suite, verificar que:

\begin{enumerate}
    \item \textbf{Todos los nodos tienen la misma imagen del backend.} Comparar con:
    \begin{lstlisting}[language=bash]
for ip in <ip-master> <ip-worker1> <ip-worker2>; do
  ssh student1@$ip "sudo k3s ctr images ls | grep e-converse-backend | awk '{print \$3}'"
done
    \end{lstlisting}
    Los SHAs deben ser idénticos. Si no, exportar desde un nodo correcto e importar en el resto (ver Sección 4 del \texttt{REPORTE\_PROYECTO.md}).

    \item \textbf{El clúster está en estado limpio:} 3 nodos en \texttt{Ready}, deployments backend/frontend en \texttt{Running}.

    \item \textbf{El endpoint responde 200 antes de empezar:}
    \begin{lstlisting}[language=bash]
curl -s -o /dev/null -w "%{http_code}\n" http://10.6.101.163:30080/api/categorias/list
    \end{lstlisting}
\end{enumerate}

% ======================================================
\section{Resumen de Resultados}
% ======================================================

Para el detalle completo (gráficas comparativas, heatmaps, curvas de escalado, análisis cuantitativo) ver el notebook. Aquí solo el resumen ejecutivo:

\begin{center}
\small
\begin{tabular}{lccccc}
\toprule
\textbf{Escenario} & \textbf{Nodos} & \textbf{Réplicas} & \textbf{Throughput Cat (req/s)} & \textbf{Latencia Cat (ms)} & \textbf{Error \%} \\
\midrule
1N-1R & 1 & 1 & 1.37 & 398 & 0.0 \\
1N-2R & 1 & 2 & 1.18 & 593 & 0.0 \\
1N-3R & 1 & 3 & 0.73 & 1{,}076 & 0.0 \\
2N-1R & 2 & 1 & 1.83 & 396 & 0.0 \\
2N-2R & 2 & 2 & 2.15 & 318 & 0.0 \\
2N-3R & 2 & 3 & 2.60 & 309 & 0.0 \\
3N-1R & 3 & 1 & 1.77 & 356 & 0.0 \\
\textbf{3N-2R} & \textbf{3} & \textbf{2} & \textbf{4.27} & \textbf{294} & \textbf{0.0} \\
3N-3R & 3 & 3 & 3.27 & 289 & 0.0 \\
\bottomrule
\end{tabular}
\end{center}

\textbf{Configuración óptima:} 3 nodos $\times$ 2 réplicas (4.27 req/s, 294 ms de latencia, 0\% errores).

\textbf{Conclusiones principales:}
\begin{enumerate}
    \item El escalado horizontal funciona, pero tiene un techo definido por la base de datos. Pasar de 2 a 3 réplicas con 3 nodos degrada el throughput ($-23$\%) por contención de connection pool en MongoDB Atlas.
    \item Concentrar la carga en un solo nodo (1N-3R) entrega menos de $1/5$ del throughput del mejor escenario.
    \item El tamaño del payload (3.5 MB del endpoint de productos) domina sobre la topología: la latencia es 10--15$\times$ la del endpoint ligero en \emph{todos} los escenarios.
\end{enumerate}

% ======================================================
\section{Enlaces y Recursos}
% ======================================================

\begin{itemize}
    \item \textbf{Repositorio principal:} \\
    \url{https://github.com/JuanSe2731/Proyecto-Spring-Boot-E-converse_vers2}

    \item \textbf{Notebook de análisis (en el repositorio):} \\
    \texttt{Analisis\_Rendimiento.ipynb}

    \item \textbf{Diario de problemas y soluciones:} \\
    \texttt{REPORTE\_PROYECTO.md}

    \item \textbf{Guía para video explicativo:} \\
    \texttt{GUIA\_VIDEO.md}

    \item \textbf{Documentación Apache JMeter:} \\
    \url{https://jmeter.apache.org/usermanual/get-started.html}

    \item \textbf{Documentación K3s:} \\
    \url{https://docs.k3s.io/}

    \item \textbf{MongoDB Atlas:} \\
    \url{https://www.mongodb.com/atlas}
\end{itemize}

\vspace{1cm}
\hrule
\vspace{0.3cm}
\begin{center}
\small\textit{Universidad Industrial de Santander --- Ingeniería de Software III --- Mayo 2026}
\end{center}

\end{document}
